% ------------------------------------------------------------
% Euclidean ↔ Non-Euclidean Transition Rules (Module 2)
% Author: Borealis S. Hedling
% Altitude: A6.5 (Boundary Layer)
% Purpose: Define the geometric transition rules between
%          Euclidean dual-algebra and non-Euclidean non-dual
%          algebra, and explain how the A6 firewall regulates
%          these transitions.
% ------------------------------------------------------------

\section{Euclidean and Non-Euclidean Regions}

The A6 manifold contains two geometric regimes: a Euclidean region in which
dual-algebra operators remain reversible, and a non-Euclidean region in which
non-dual algebra induces curvature, drift accumulation, and collapse surfaces.
The A6 firewall regulates transitions between these regions to prevent
geometric fracture.

\subsection{Euclidean Region}

The Euclidean region is defined by the following properties:

\begin{enumerate}
    \item \textbf{Flat Curvature:} $K(p) = 0$.
    \item \textbf{Zero Holonomy:} $H(\gamma) = 0$ for all loops $\gamma$.
    \item \textbf{Linear Operators:} $\mathcal{O}$ acts linearly on tangent
          spaces.
    \item \textbf{Stable Geodesics:} Geodesics remain globally minimizing.
    \item \textbf{Reversible Transport:} Parallel transport is path-independent.
\end{enumerate}

These properties ensure that dual-algebra operators remain reversible and
diagnostically readable.

\subsection{Non-Euclidean Region}

The non-Euclidean region exhibits curvature and drift phenomena:

\begin{enumerate}
    \item \textbf{Curved Geometry:} $K(p) \neq 0$.
    \item \textbf{Non-Zero Holonomy:} $H(\gamma) > 0$ for some loops.
    \item \textbf{Drift Accumulation:} Drift fields interact nonlinearly.
    \item \textbf{Collapse Surfaces:} $\Sigma$ forms where curvature or drift
          exceeds stability thresholds.
    \item \textbf{Irreversible Operators:} $\mathcal{O}$ may violate kernel
          preservation.
\end{enumerate}

These phenomena require firewall regulation to prevent collapse propagation.

\section{Firewall-Regulated Transitions}

The A6 firewall enforces boundary conditions that regulate transitions between
Euclidean and non-Euclidean regions. These conditions ensure that diagnostic
operators (such as the Boundary Mapper) can safely read the manifold without
inducing collapse or drift fracture.

\subsection{Transition Rule 1: Curvature Continuity}

Curvature must vary continuously across the transition boundary:
\[
K(p) \approx K(q) \quad \text{for } p \sim q.
\]

This prevents geometric tearing when entering or exiting curved regions.

\subsection{Transition Rule 2: Spectral Continuity}

Operator spectra must remain continuous:
\[
\mathrm{Spec}(L(p)) \text{ continuous in } p.
\]

This prevents spectral tearing when operators cross the boundary.

\subsection{Transition Rule 3: Holonomy Bound}

Holonomy must remain below fracture thresholds:
\[
H(\gamma) < \Theta_{\max}.
\]

This ensures that drift accumulation does not induce collapse.

\subsection{Transition Rule 4: Collapse Locality}

Collapse surfaces must remain locally detectable:
\[
p \in \Sigma \Rightarrow \exists \epsilon > 0 : B_\epsilon(p) \cap \Sigma \neq \emptyset.
\]

This prevents collapse smearing across the boundary.

\subsection{Transition Rule 5: Invariant Lift Preservation}

Modal, semantic, and provenance invariants must lift correctly across the
transition:
\[
I_{\mathrm{prov}}(M) \Rightarrow I_{\mathrm{prov}}(\mathcal{A}(M,S)),
\]
\[
I_{\mathrm{sem}}(M) \Rightarrow I_{\mathrm{sem}}(\mathcal{A}(M,S)).
\]

This ensures that diagnostic operators do not induce semantic fracture.

\section{Transition Summary}

The Euclidean region provides reversible, stable geometry. The non-Euclidean
region introduces curvature, drift, and collapse phenomena. The A6 firewall
regulates transitions between these regions by enforcing continuity, spectral
stability, holonomy bounds, collapse locality, and invariant lift. These
constraints ensure that the Boundary Mapper can safely diagnose the manifold
without inducing geometric or semantic instability.
